%=======================================================================
%  SNIPPET GALLERY & FORMATTING EXAMPLES
%  Copy any snippet below into your weekly sections/03_results.tex as needed.
%=======================================================================

%-------------------------------------------------- 1. TikZ Diagram ----
\subsection{System Architecture Diagram}
\begin{figure}[H]
\centering
\begin{tikzpicture}[
  node distance=8mm and 9mm,
  box/.style={draw=accent!60,fill=accent!7,rounded corners=2pt,
              minimum height=9mm,minimum width=23mm,align=center,
              font=\small,inner sep=3pt},
  store/.style={box,fill=ok!8,draw=ok!55},
  outbox/.style={box,fill=warn!10,draw=warn!60},
  arr/.style={-{Stealth[length=2mm]},draw=neutral!80}]

  \node[box] (src)   {Dataset\\stream};
  \node[box,right=of src] (drift) {Drift\\injector};
  \node[box,right=of drift] (algo) {Estimator};
  \node[store,right=of algo] (met) {Metric\\registry};
  \node[outbox,below=of met] (rep) {Report};

  \draw[arr] (src) -- (drift);
  \draw[arr] (drift) -- (algo);
  \draw[arr] (algo) -- (met);
  \draw[arr] (met) -- (rep);
\end{tikzpicture}
\caption{Proposed system pipeline.}
\label{fig:pipeline_example}
\end{figure}
\takeaway{Component decoupling ensures reproducible benchmarking.}

%-------------------------------------------------- 2. Single Figure ---
\subsection{Single Figure}
\begin{figure}[H]
\centering
\includegraphics[width=0.75\linewidth]{figure_name.png}
\caption{Experiment visualization.}
\label{fig:single_example}
\end{figure}
\takeaway{Summary of observations from the figure above.}

%-------------------------------------------------- 3. Side-by-side Subfigures
\subsection{Side-by-Side Subfigures}
\begin{figure}[H]
\centering
\begin{subfigure}[b]{0.48\linewidth}
  \centering
  \includegraphics[width=\linewidth]{before.png}
  \caption{Before treatment}
  \label{fig:before_example}
\end{subfigure}\hfill
\begin{subfigure}[b]{0.48\linewidth}
  \centering
  \includegraphics[width=\linewidth]{after.png}
  \caption{After treatment}
  \label{fig:after_example}
\end{subfigure}
\caption{Comparison before and after treatment.}
\label{fig:compare_example}
\end{figure}
\takeaway{Key difference observed between the two conditions.}

%-------------------------------------------------- 4. Figure Next to Table
\subsection{Figure Next to a Table}
\begin{figure}[H]
\centering
\begin{minipage}[t]{0.46\linewidth}
  \centering
  \vspace{0pt}
  \includegraphics[width=\linewidth]{memory.png}
  \captionof{figure}{Memory profile over time.}
  \label{fig:mem_example}
\end{minipage}\hfill
\begin{minipage}[t]{0.5\linewidth}
  \vspace{0pt}
  \centering
  \captionof{table}{Peak resource usage.}
  \label{tab:mem_example}
  \small
  \begin{tabular}{@{}lrr@{}}
  \toprule
  \textbf{Algorithm} & \textbf{Peak RAM} & \textbf{Throughput} \\
  \midrule
  Model A & 41\,MB & 12\,400 pts/s \\
  Model B & 78\,MB &  6\,900 pts/s \\
  \bottomrule
  \end{tabular}
\end{minipage}
\end{figure}
\takeaway{Comparison between memory footprint and execution speed.}

%-------------------------------------------------- 5. Text Wrapped Around Figure
\subsection{Text Wrapped Around Figure}
\begin{wrapfigure}{r}{0.35\linewidth}
  \vspace{-8pt}
  \centering
  \includegraphics[width=\linewidth]{thumbnail.png}
  \caption{Radius growth.}
  \label{fig:wrap_example}
  \vspace{-6pt}
\end{wrapfigure}
Insert descriptive text here explaining the wrapped figure. Keep at least 4 to 5 lines
of text next to a wrapped figure to ensure proper alignment.

%-------------------------------------------------- 6. Full Results Table
\subsection{Results Table}
\begin{table}[H]
\centering
\caption{Benchmark results. Best values in \textbf{bold}.}
\label{tab:results_example}
\small
\begin{tabular}{@{}l cc cc@{}}
\toprule
& \multicolumn{2}{c}{\textbf{Dataset A}} & \multicolumn{2}{c}{\textbf{Dataset B}} \\
\cmidrule(lr){2-3}\cmidrule(lr){4-5}
\textbf{Model} & Accuracy & F1-Score & Accuracy & F1-Score \\
\midrule
Baseline & 0.71 & 0.55 & 0.63 & 0.48 \\
Ours     & \textbf{0.84} & \textbf{0.79} & \textbf{0.78} & \textbf{0.67} \\
\bottomrule
\end{tabular}
\end{table}
\takeaway{Our method consistently outperforms the baseline across both datasets.}

%-------------------------------------------------- 7. Bar Chart (PGFPlots)
\subsection{Performance Bar Chart}
\begin{figure}[H]
\centering
\begin{tikzpicture}
\begin{axis}[
  ybar, width=0.86\linewidth, height=5.4cm,
  bar width=11pt, enlarge x limits=0.28,
  symbolic x coords={ModelA,ModelB,ModelC},
  xtick=data, ymin=0, ymax=1.0,
  ylabel={Score}, ylabel style={font=\small},
  tick label style={font=\small},
  legend style={at={(0.5,1.03)},anchor=south,legend columns=-1,
                font=\small,draw=rulegray},
  nodes near coords, every node near coord/.append style={font=\tiny},
  grid=major, grid style={rulegray,dashed}]
  \addplot[fill=accent!65,draw=accent!80] coordinates {(ModelA,0.71) (ModelB,0.64) (ModelC,0.85)};
  \addplot[fill=ok!55,draw=ok!75] coordinates {(ModelA,0.55) (ModelB,0.49) (ModelC,0.78)};
  \legend{Metric 1, Metric 2}
\end{axis}
\end{tikzpicture}
\caption{Performance comparison across models.}
\label{fig:chart_example}
\end{figure}
\takeaway{Model C achieves superior performance on both metrics.}

%-------------------------------------------------- 8. Code Listing
\subsection{Code Snippet}
\begin{lstlisting}[style=py,caption={Evaluation loop implementation.},label={lst:code_example}]
for batch in dataloader:
    pred = model(batch)
    loss = criterion(pred, batch.target)
    loss.backward()
    optimizer.step()
\end{lstlisting}
\takeaway{Brief explanation of key implementation details in the snippet above.}
